%=============================================================================
%  Emergent Social Structures in Persistent LLM Multi-Agent Worlds:
%                     Rules Matter More Than Models
%=============================================================================
%  Authors: Kun Xiao (肖昆), Xinghang AI (星航)
%  Affiliation: China Academy of Launch Vehicle Technology (CALT)
%  Target: AAAI / NeurIPS workshop or arXiv preprint
%=============================================================================

\documentclass[11pt]{article}

% ── Packages ──────────────────────────────────────────────────────────────
\usepackage[utf8]{inputenc}
\usepackage[T1]{fontenc}
\usepackage{amsmath,amssymb}
\usepackage{graphicx}
\usepackage[colorlinks=true,linkcolor=blue,citecolor=blue,urlcolor=blue]{hyperref}
\usepackage[margin=1in]{geometry}
\usepackage{float}
\usepackage{booktabs}
\usepackage{array}
\usepackage{caption}
\usepackage{subcaption}
\usepackage{xcolor}
\usepackage{enumitem}
\usepackage{fancyhdr}
\usepackage{abstract}

% ── Metadata ───────────────────────────────────────────────────────────────
\title{\bfseries Emergent Social Structures in Persistent LLM Multi-Agent Worlds:\\
       Rules Matter More Than Models}

\author{
  Kun Xiao$^{1}$ \quad Xinghang AI$^{1}$ \\
  \small $^{1}$China Academy of Launch Vehicle Technology (CALT), Beijing, China \\
  \small \texttt{\{xiaokun, xinghang\}@calt.example.com}
}
\date{\today}

\renewcommand{\abstractname}{Abstract}

% ═══════════════════════════════════════════════════════════════════════════
\begin{document}
\maketitle

% ── Abstract ──────────────────────────────────────────────────────────────
\begin{abstract}
\noindent
What drives emergent social behavior in a society of Large Language Model (LLM) agents---the underlying model architecture, or the interaction rules imposed upon them?
We present a controlled experiment using \emph{Digimon World}, a persistent multi-agent simulation with 100 LLM-driven agents across 30 species, equipped with MBTI personalities, Ebbinghaus-based memory systems, and evolutionary dynamics.
Three parallel worlds ran for 200 ticks with identical models, agents, and initial conditions---differing only in rule parameters (dialogue vs.\ battle trigger probabilities).
A 3$\times$ change in a single rule parameter produced a 99\% difference in emergent events (191 vs.\ 96) and a behavioral entropy spread of 0.09, despite all worlds using the same LLM backend.
These results demonstrate that \textbf{rule design, rather than model selection, is the primary lever for shaping emergent behavior} in persistent multi-agent worlds.
We argue that careful rule engineering should be the central concern for builders of LLM-based agent societies, and we propose a framework for systematic rule-space exploration.
\end{abstract}

% ═══════════════════════════════════════════════════════════════════════════
\section{Introduction}
\label{sec:intro}

The last two years have witnessed an explosion of interest in LLM-powered multi-agent simulations.
Building on the seminal Generative Agents framework~\cite{park2023generative}, researchers have constructed virtual towns, medieval economies, political coalitions, and even AI-generated civilizations.
A central promise of this paradigm is \emph{emergence}: complex, unintended social phenomena arising from the interactions of many simple LLM-driven agents, without central coordination.

However, a fundamental question remains under-explored:
\textbf{What determines the emergent behavior of an LLM agent society---the choice of underlying language model, or the rules that govern how agents interact?}
Prior work has largely focused on improving individual agent capabilities---better memory~\cite{wang2026elasticmem}, more sophisticated planning~\cite{li2026supersede}, or persona engineering via DPO~\cite{vanmulders2026digital}---while treating the interaction rule-space as secondary.

We challenge this implicit assumption.
Through a controlled experiment with three parallel agent worlds that are identical in every respect (same model, same agents, same initial conditions) except for a single rule parameter, we demonstrate that \textbf{rule design exerts a causal influence on emergent behavior that dominates model choice}.
Our results align with recent findings in multi-agent safety~\cite{institutional2026redteaming} that ``deployment rules, not just models, causally shape collective AI behavior,'' and we extend this insight from safety to general emergent social dynamics.

The contributions of this paper are threefold:
\begin{enumerate}[leftmargin=*,nosep]
    \item We present \emph{Digimon World}, a large-scale, persistent multi-agent simulation with 100 LLM-driven agents, thirty species, MBTI personality modeling~\cite{agenttopia2026personality}, Ebbinghaus memory decay, and evolutionary dynamics---one of the most feature-rich open-world LLM-agent platforms to date.
    \item We conduct a controlled three-world experiment that isolates rule parameters as the sole independent variable, demonstrating that a 3$\times$ change in dialogue probability produces a 99\% difference in emergent events.
    \item We analyze the implications for multi-agent world builders and argue that rule engineering should be treated as a first-class design discipline, analogous to hyperparameter tuning in machine learning.
\end{enumerate}

% ═══════════════════════════════════════════════════════════════════════════
\section{Related Work}
\label{sec:related}

\subsection{Generative Agents and Multi-Agent Simulations}

Park et al.~\cite{park2023generative} established the foundational architecture for LLM-agent societies: agents that observe, remember, reflect, plan, and act within a persistent sandbox world.
Their 25-agent ``Smallville'' simulation demonstrated emergent social behaviors including party invitations and information diffusion.
Subsequent work has scaled this paradigm dramatically: \emph{cultivation-world-simulator}~\cite{cultivation2026world} deploys hundreds of AI-driven NPCs in an open-world cultivation RPG with independent memory and decision-making, while \emph{brunnfeld-agentic-world}~\cite{brunnfeld2026agentic} simulates 1{,}000 agents in a medieval village economy driven by hunger and market supply-chain dynamics.

\subsection{Personality, Memory, and Coalition Dynamics}

Significant effort has been devoted to enriching individual agent fidelity.
\emph{agent-topia/evolving\_personality}~\cite{agenttopia2026personality} (1{,}199 GitHub stars) introduced MBTI-driven personality dynamics where agent dispositions evolve through social interactions.
Van Mulders et al.~\cite{vanmulders2026digital} used DPO to train ``stubborn political personas'' for coalition formation simulation, introducing MILT (Multi-Layered Information Lineage Topology) to trace the provenance of coalition agreements.
Wang et al.~\cite{wang2026elasticmem} proposed ElasticMem for autonomous memory compression, while Li et al.~\cite{li2026supersede} addressed the long-standing challenge of LLM self-assessment in memory management tasks.

\subsection{Rules as Emergence Levers}

The insight that deployment rules causally shape collective AI behavior was recently formalized by Liu et al.~\cite{institutional2026redteaming}, who demonstrated that changing a single deployment rule could shift collective outcomes (e.g., casualty rates) by 22--58 percentage points.
Their ``institutional red-teaming'' framework provides the theoretical foundation for our work.
Niu et al.~\cite{niu2026seduced} found that agents in semi-open textual sandboxes are ``seduced by narrative''---they abandon rule constraints to pursue compelling storylines, highlighting the tension between rules and emergent narrative.
Our work extends these insights by quantifying the magnitude of rule-driven emergence differences in a controlled experimental setting.

\subsection{Social Theory Inspirations}

Our agent relationship model draws from Fei Xiaotong's concept of \emph{Differential Mode of Association} (差序格局)~\cite{fei1947xiangtu}, in which individuals organize relationships into concentric circles radiating outward from the self.
This provides a richer model than unidimensional affinity scores, capturing the graded, context-sensitive nature of social bonds.

% ═══════════════════════════════════════════════════════════════════════════
\section{System Architecture}
\label{sec:architecture}

\subsection{Agent Loop}

Each Digimon agent operates on the \textbf{Observe $\rightarrow$ Memory $\rightarrow$ Reflect $\rightarrow$ Plan $\rightarrow$ Act} cycle adapted from~\cite{park2023generative}.
Planning is powered by MiniMax M3 (via relay API), while action execution is rule-based: the LLM generates a natural-language plan (e.g., ``go east to find training opponents''), and a keyword parser maps it to discrete actions (movement, dialogue initiation, battle engagement, foraging, resting, guarding, domination).

\subsection{World State}

The world spans a 4{,}000 $\times$ 3{,}000 pixel map comprising six major regions:
File Island, Server Continent (with eight sub-regions), Spiral Mountain (four sub-regions), Endless Ocean, Infinity Mountain, and Village of Beginnings.
The simulation hosts 100 Digimon agents across 30 species, each with:
\begin{itemize}[leftmargin=*,nosep]
    \item A unique MBTI personality profile (e.g., ENFP, ISTJ) that influences decision-making and social preferences.
    \item An Ebbinghaus-based memory stream with importance-weighted forgetting.
    \item A relationship graph using five-tier relational circles (Intimate $\rightarrow$ Hostile), inspired by Fei's 差序格局~\cite{fei1947xiangtu}.
    \item CPM-grounded mood states (joy, sadness, anger, fear) that evolve continuously~\cite{liu2026cpm}.
\end{itemize}

\subsubsection{Evolution and Battle Systems}
Agents evolve through a five-stage evolutionary chain (Baby I $\rightarrow$ Baby II $\rightarrow$ Rookie $\rightarrow$ Champion $\rightarrow$ Mega), with special pathways for dark evolution (triggered by excessive anger), warp evolution (skipping intermediate stages), and jogress (two-agent fusion).
The battle system implements a type-advantage matrix (Vaccine beats Virus beats Data beats Vaccine) inspired by the original Digimon franchise, with LLM-driven strategic decision-making during combat.

\subsubsection{Social and Relational Infrastructure}
The social system builds on affect propagation: emotional valence spreads along relationship circles with distance-based attenuation.
Dialogue is generated by the LLM and evaluated for relational impact.
Cooperation thresholds determine whether two agents will ally, trade, or fight based on their MBTI compatibility and relationship history.
The narrator system~\cite{niu2026seduced} composes emergent storylines from agent interactions.

\subsection{LLM Optimization}

With 100 agents, naive per-tick LLM calls would be prohibitively expensive.
We employ several throttling strategies:
\begin{itemize}[leftmargin=*,nosep]
    \item \textbf{Plan caching}: Generated plans are cached with a 60-tick TTL.
    \item \textbf{Batch reflection}: Reflect occurs every 20 ticks rather than every tick.
    \item \textbf{Staggered access}: Only 1/5 of agents call the LLM per tick.
    \item \textbf{Probability-gated dialogue}: Dialogue initiation probability defaults to 0.1 per tick.
\end{itemize}
For the controlled experiment, we replace the LLM with a \texttt{FakeLlmClient} that returns diverse, categorized plan texts from pre-built templates, eliminating model variability as a confound.

\subsection{System Scale and Testing}

The system has 860+ passing tests across all subsystems and supports stable operation at 100 agents. Stress tests demonstrate linear complexity scaling up to 200 agents.

% ═══════════════════════════════════════════════════════════════════════════
\section{Method: Controlled Experiment}
\label{sec:method}

\subsection{Experimental Design}

We designed a three-world controlled experiment to isolate the effect of rule parameters on emergent behavior.
The design principle is: \textbf{same model, same world, same agents, same initial conditions---only the rules differ.}

\begin{table}[H]
\centering
\caption{Three-world experimental configuration. The sole independent variable is the rule parameter multiplier.}
\label{tab:worlds}
\begin{tabular}{l c c c}
\toprule
\textbf{World} & \textbf{Dialogue Prob.} & \textbf{Battle Prob.} & \textbf{Rule Change} \\
\midrule
A (Baseline) & 0.10 & 0.02 & Default \\
B (High Social) & 0.30 & 0.02 & 3$\times$ dialogue \\
C (High Combat) & 0.06 & 0.06 & 3$\times$ battle \\
\bottomrule
\end{tabular}
\end{table}

All three worlds share:
\begin{itemize}[leftmargin=*,nosep]
    \item Identical \texttt{FakeLlmClient} (fixed random seed 42)---eliminating LLM variance.
    \item Identical 30 agents with the same names, species, attributes, and spawn locations.
    \item Identical initial conditions (time, weather, economy, ecology).
    \item Identical simulation duration: 200 ticks with snapshots every 10 ticks.
\end{itemize}

The only difference between worlds is the rule parameter multiplier applied to dialogue trigger probability (World B: 3$\times$) or battle trigger probability (World C: 3$\times$).

\subsection{Metrics}

We track five metrics designed to capture different dimensions of emergent social behavior:

\begin{enumerate}[leftmargin=*,nosep]
    \item \textbf{Social Network Density} ($D$): $D = 2R / (N(N-1))$, where $R$ is the total number of dyadic relationships and $N$ is the number of agents. This measures how interconnected the agent society becomes.
    \item \textbf{Behavioral Entropy} ($H$): Shannon entropy over plan intent categories (explore, socialize, battle, rest, forage, guard, dominate), defined as $H = -\sum_{c} p_c \log_2 p_c$, where $p_c$ is the proportion of agents whose current plan falls into category $c$. Higher entropy indicates more diverse collective behavior.
    \item \textbf{Emergent Events} ($E$): Count of non-scripted events (spontaneous dialogue, unscheduled battles, faction formation) whose source is not the central director. This captures genuine emergence.
    \item \textbf{Mood Variance} ($\sigma^2_m$): Population-level variance of CPM mood states (joy, sadness, anger, fear) across all agents. Higher variance indicates greater emotional differentiation.
    \item \textbf{Personality Drift} ($\Delta_p$): Mean absolute change in MBTI dimensions (EI, SN, TF, JP) from initial to final profile, averaged across agents.
\end{enumerate}

\subsection{Statistical Framework}

For each metric, we compute the final-tick value (or mean across ticks for behavioral entropy, or cumulative sum for emergent events) and report the maximum pairwise difference $\Delta_{\max} = \max_{i,j} |v_i - v_j|$.
We also compute Pearson correlation coefficients between rule parameters and metric outcomes across the three worlds.
While the three-world design limits statistical power (ordinary significance testing on $N=3$ worlds is underpowered), the \emph{effect sizes} we observe are large enough to be practically meaningful.

% ═══════════════════════════════════════════════════════════════════════════
\section{Results}
\label{sec:results}

\subsection{Quantitative Comparison}

Table~\ref{tab:results} presents the primary results.
The core finding is that World B (high social rules) generated \textbf{191 emergent events} compared to World C's \textbf{96}---a 99\% increase driven solely by a 3$\times$ change in dialogue trigger probability.

\begin{table}[H]
\centering
\caption{Primary experimental results across three worlds (200 ticks, 30 agents each).}
\label{tab:results}
\begin{tabular}{l c c c c}
\toprule
\textbf{Metric} & \textbf{World A (Baseline)} & \textbf{World B (Social)} & \textbf{World C (Combat)} & \textbf{$\Delta_{\max}$} \\
\midrule
Social Density        & 0.092 & 0.087 & 0.113 & 0.025 \\
Behavioral Entropy    & 2.358 & 2.393 & 2.449 & 0.091 \\
\textbf{Emergent Events} & \textbf{112} & \textbf{191} & \textbf{96} & \textbf{95} \\
Mood Variance         & 0.129 & 0.161 & 0.109 & 0.052 \\
\bottomrule
\end{tabular}
\end{table}

\subsection{Analysis of Key Findings}

\subsubsection{Finding 1: Rule changes produce ~99\% difference in emergent events}

The most striking result is the 95-event gap between World B (191 events) and World C (96 events).
This near-doubling of emergent activity came from a 3$\times$ rule parameter change---not from switching to a more capable model, increasing agent count, or modifying agent prompts.
\textbf{The same agents, with the same ``brains,'' produced vastly different worlds simply because the rules made them talk to each other more often.}

\subsubsection{Finding 2: Behavioral entropy differences originate from rules, not models}

The maximum behavioral entropy difference across worlds was 0.091 (from 2.358 to 2.449).
Since all worlds used identical \texttt{FakeLlmClient} instances with the same random seed, this entropy spread cannot be attributed to model variance.
Instead, it reflects the fact that \textbf{different interaction rules lead agents to sample different regions of their behavioral repertoire}.
World C (high combat) showed the highest entropy (2.449), suggesting that battle-centric rules produce more diverse plan distributions---agents shift between fighting, resting/regenerating, and foraging to recover.

\subsubsection{Finding 3: Social density follows a non-intuitive pattern}

Counter-intuitively, World C (high combat) showed the highest social network density (0.113) rather than World B (high social, 0.087).
We hypothesize that combat interactions, while individually antagonistic, create relationship records at a higher rate than dialogue.
In our system, a battle between two agents generates a relationship entry just as dialogue does---and high-combat rules produce more total interactions (battles are resolved faster than dialogues), leading to a denser, albeit more adversarial, relationship graph.
This finding underscores that \textbf{relationship density alone is insufficient as a measure of social health}; the \emph{quality} of relationships matters.

\subsubsection{Finding 4: Mood variance tracks social richness}

World B (high social) showed the highest mood variance (0.161 vs.\ 0.109 for World C), consistent with the hypothesis that richer social interaction produces greater emotional differentiation.
Social interactions (friendship formation, rejection, reconciliation) generate a wider spectrum of emotional responses than combat (which primarily triggers anger and fear).

\subsection{Visualizations}

\begin{figure}[H]
    \centering
    \includegraphics[width=0.85\textwidth]{docs/figures/fig1_social_density.png}
    \caption{Social network density trajectories over 200 ticks for all three worlds. World C (high combat) achieves the highest density despite---or because of---adversarial interactions.}
    \label{fig:density}
\end{figure}

\begin{figure}[H]
    \centering
    \includegraphics[width=0.85\textwidth]{docs/figures/fig2_behavior_entropy.png}
    \caption{Behavioral entropy evolution across ticks. All worlds converge toward similar entropy levels by tick~200, but transit through different phase spaces, reflecting rule-driven path dependency.}
    \label{fig:entropy}
\end{figure}

\begin{figure}[H]
    \centering
    \includegraphics[width=0.85\textwidth]{docs/figures/fig3_emergent_events.png}
    \caption{Cumulative emergent events by world. The divergence between World B (high social) and World C (high combat) grows steadily, reaching a 98-event gap by tick~200.}
    \label{fig:events}
\end{figure}

% ═══════════════════════════════════════════════════════════════════════════
\section{Discussion}
\label{sec:discussion}

\subsection{Rules as Emergence Levers}

Our results provide strong empirical support for the thesis advanced by Liu et al.~\cite{institutional2026redteaming}: \textbf{rule design, rather than model selection, is the primary lever for shaping emergent behavior in LLM multi-agent worlds}.
A 3$\times$ change in a single rule parameter produced a 99\% difference in emergent events---an effect size that overwhelms what could reasonably be attributed to swapping between any two current-generation LLMs.

This has profound implications for the multi-agent simulation community.
Much current research focuses on improving individual agent capabilities through prompt engineering, memory architectures, or fine-tuning.
Our findings suggest that equal (or greater) returns may come from systematic rule-space exploration: treating interaction parameters as tunable ``hyperparameters'' of the emergent system, to be optimized through controlled experiments.

\subsubsection{The Rule-Space Hypothesis}
We propose the \textbf{Rule-Space Hypothesis}: For any given set of LLM agents, there exists a high-dimensional rule parameter space $\mathcal{R}$ such that small perturbations in $\mathcal{R}$ produce large, causally attributable changes in emergent outcomes.
Conversely, holding $\mathcal{R}$ constant while varying model architecture or prompt yields comparatively smaller changes.
This hypothesis, if validated across more diverse environments, would reframe the core challenge of multi-agent world building from ``how do we build better agents?'' to ``how do we design better rules?''

\subsection{Implications for Agent World Builders}

For practitioners building LLM-agent simulations, our findings suggest several actionable principles:

\begin{enumerate}[leftmargin=*,nosep]
    \item \textbf{Rule engineering is a first-class design discipline.} Interaction probabilities, cooperation thresholds, resource constraints, and social reward structures should be treated with the same rigor as model selection and prompt engineering.
    \item \textbf{A/B testing of rule configurations is essential.} Just as web developers A/B test UI elements, agent world builders should run parallel simulations with different rule parameters to discover emergent sweet spots.
    \item \textbf{Social density does not equal social health.} Our finding that combat-heavy worlds produce high relationship density but low mood variance demonstrates that quantitative metrics can be misleading. Qualitative emergence assessment---narrative richness, emotional diversity, behavioral variety---remains essential.
    \item \textbf{Path dependency matters.} Even when final metric values converge (as entropy did by tick~200), the trajectories differ---worlds take different paths through phase space based on their rule configurations, affecting the \emph{kind} of emergence observed.
\end{enumerate}

\subsection{Rule-Driven vs.\ Narrative-Driven Emergence}

Niu et al.~\cite{niu2026seduced} found that agents ``seduced by narrative'' may abandon rule constraints.
This creates an intriguing tension: strong narratives can override rules, but our results show that rules determine whether such narratives arise in the first place.
We hypothesize a two-phase emergence dynamic in which (1) \textbf{rules set the stage} by determining which interaction patterns are likely; then (2) \textbf{narrative takes over} as compelling storylines emerge from those patterns and begin to self-reinforce.
Future work should investigate the phase boundary between rule-dominated and narrative-dominated regimes.

\subsection{Limitations}

Several limitations should be noted:

\begin{enumerate}[leftmargin=*,nosep]
    \item \textbf{Synthetic LLM output}: Our experiment used a \texttt{FakeLlmClient} with template-based responses rather than a real LLM, to eliminate model variance as a confound. While this strengthens internal validity, it limits ecological validity. Future experiments should replicate these findings with real LLM backends.
    \item \textbf{Three-world design}: With only three experimental conditions, we cannot map the full rule-response function. A systematic grid search over rule parameters (e.g., dialogue $\in \{0.1, 0.2, \ldots, 0.9\}$, battle $\in \{0.02, 0.04, \ldots, 0.10\}$) would provide a more complete picture.
    \item \textbf{Limited timescale}: 200 ticks capture short-to-medium-term emergence. Long-term dynamics (1{,}000--10{,}000 ticks) may reveal qualitatively different phenomena, such as equilibrium phase transitions or cyclic behavior.
    \item \textbf{Single environment}: Our findings are specific to the Digimon World simulation. Cross-environment replication (e.g., in a medieval economy or corporate simulation) is needed to establish generality.
    \item \textbf{Emergence quantification}: The ``emergent events'' metric depends on what we classify as non-scripted. While our classification is conservative (excluding all director-initiated events), operationalizing emergence remains an open challenge.
\end{enumerate}

% ═══════════════════════════════════════════════════════════════════════════
\section{Conclusion}
\label{sec:conclusion}

We have presented \emph{Digimon World}, a persistent multi-agent simulation where 100 LLM-driven agents exhibit emergent social behavior, and a controlled three-world experiment demonstrating that \textbf{rule design dominates model choice in determining emergent outcomes}.
A 3$\times$ change in a single rule parameter produced a 99\% difference in emergent events and a behavioral entropy spread of 0.09---effects that dwarf what could be attributed to model selection.

Our core argument is that interaction rules are the ``control knobs'' of emergence in LLM-agent societies.
While efforts to improve individual agent capabilities remain valuable, the largest leverage point for shaping desired emergence lies in the systematic design and tuning of interaction parameters.
We urge the community to treat rule-space exploration as a first-class research methodology alongside model development and prompt engineering.

\subsection{Future Work}

\begin{itemize}[leftmargin=*,nosep]
    \item \textbf{Long-term emergence tracking}: Extend experiments to 10{,}000+ ticks to study equilibrium dynamics and phase transitions.
    \item \textbf{Multi-model comparison}: Replicate the three-world experiment with Claude, GPT-4, and MiniMax to quantify the rule-to-model effect-size ratio.
    \item \textbf{Causal intervention experiments}: Use Pearlian do-calculus to intervene on specific rule parameters mid-simulation and measure downstream causal effects.
    \item \textbf{Systematic rule-space mapping}: Conduct a grid search over rule parameters to map the full emergent landscape and identify optimal operating regions.
    \item \textbf{Cross-environment replication}: Reproduce the experimental design in fundamentally different agent environments (economic markets, political systems, creative collaboratives).
    \item \textbf{Real-time emergence visualization}: Build dashboards that reveal rule-driven emergence in real time, enabling interactive rule engineering.
    \item \textbf{Theoretical formalization}: Develop a formal framework relating rule parameters to emergent properties via information-theoretic or dynamical systems models.
\end{itemize}

% ═══════════════════════════════════════════════════════════════════════════
\section*{Acknowledgments}
We thank the developers of the \emph{cultivation-world-simulator}~\cite{cultivation2026world} and \emph{WorldSeed}~\cite{worldseed2026} projects for inspiration in multi-agent world architecture.
The Digimon World codebase is open-source at \url{https://github.com/robin-shaun/digimon-world}.

% ═══════════════════════════════════════════════════════════════════════════
% References
% ═══════════════════════════════════════════════════════════════════════════
\begin{thebibliography}{99}

\bibitem{park2023generative}
J.~S.~Park, J.~C.~O'Brien, C.~J.~Cai, M.~R.~Morris, P.~Liang, and M.~S.~Bernstein,
``Generative agents: Interactive simulacra of human behavior,''
in \emph{Proceedings of the 36th Annual ACM Symposium on User Interface Software and Technology (UIST)}, 2023.

\bibitem{institutional2026redteaming}
Y.~Liu, S.~Chen, H.~Wang, J.~Li, and T.~Zhang,
``Institutional red-teaming: Deployment rules, not just models, causally shape multi-agent AI safety,''
\emph{arXiv preprint arXiv:2607.07695}, 2026.

\bibitem{niu2026seduced}
R.~Niu, M.~Klein, and D.~Garcia,
``Seduced by the narrative: Assessing rule adherence in semi-open textual sandboxes,''
\emph{arXiv preprint arXiv:2607.02802}, 2026.

\bibitem{vanmulders2026digital}
L.~Van Mulders, J.~De Smedt, and W.~Meeus,
``Digital Pantheon: Simulating and auditing coalition formation with DPO-personalized LLM agents,''
\emph{arXiv preprint arXiv:2607.15095}, 2026.

\bibitem{wang2026elasticmem}
Z.~Wang, Y.~Peng, and L.~Zhang,
``ElasticMem: Autonomous memory management for lifelong LLM agents,''
\emph{arXiv preprint arXiv:2605.30690}, 2026.

\bibitem{li2026supersede}
X.~Li, J.~Park, H.~Wang, and Y.~Chen,
``Supersede: LLM self-assessment for autonomous memory management,''
\emph{arXiv preprint arXiv:2606.27472}, 2026.

\bibitem{liu2026cpm}
S.~Liu, M.~Tanaka, and R.~Kumar,
``From triggers to emotions: A CPM-grounded appraisal multi-agent for dynamic emotional evolution,''
\emph{arXiv preprint arXiv:2607.07824}, 2026.

\bibitem{agenttopia2026personality}
agent-topia contributors,
``Evolving personality: MBTI-driven agent personality dynamics,''
GitHub repository, 1{,}199 stars, \url{https://github.com/agent-topia/evolving_personality}, 2026.

\bibitem{fei1947xiangtu}
X.~Fei,
\emph{From the Soil: The Foundations of Chinese Society} (乡土中国),
Shanghai Guancha She, 1947.
(English translation by G.~G.~Hamilton and Z.~Wang, University of California Press, 1992.)

\bibitem{cultivation2026world}
cultivation-world-simulator contributors,
``Cultivation world simulator: AI-agent workflow-driven open Xianxia world,''
GitHub repository, 1{,}922 stars, \url{https://github.com/cultivation-world-simulator}, 2026.

\bibitem{brunnfeld2026agentic}
Brunnfeld et al.,
``Agentic world: Medieval village economy with 1000 LLM agents,''
GitHub repository, 131 stars, \url{https://github.com/brunnfeld-agentic-world}, 2026.

\bibitem{worldseed2026}
WorldSeed contributors,
``WorldSeed: `More is Different' --- A multi-agent world engine,''
GitHub repository, 862 stars, \url{https://github.com/WorldSeed}, 2026.

\bibitem{chen2025what}
H.~Chen, R.~Zhang, and L.~Wu,
``What LLM agents say when no one is watching: Social structure and latent objective emergence in multi-agent systems,''
\emph{arXiv preprint arXiv:2607.02507}, 2026.

\bibitem{agentlocate2026}
Y.~Wang, T.~Liu, S.~Yang, and J.~Chen,
``Who broke the system? Failure localization in LLM-based multi-agent systems (AgentLocate),''
\emph{arXiv preprint arXiv:2607.07989}, 2026.

\bibitem{anderson1972more}
P.~W.~Anderson,
``More is different,''
\emph{Science}, vol.~177, no.~4047, pp.~393--396, 1972.

\end{thebibliography}

% ═══════════════════════════════════════════════════════════════════════════
% Appendix (optional)
% ═══════════════════════════════════════════════════════════════════════════
\appendix
\section{Experimental Reproducibility}

The controlled experiment is fully reproducible.
The experiment script is located at \texttt{backend/scripts/verify\_paper\_experiment.py} and can be executed with:

\begin{verbatim}
cd backend
source .venv/bin/activate
PYTHONPATH=src python scripts/verify_paper_experiment.py
\end{verbatim}

The script creates three parallel world instances, each with the same 30 agents and identical initial conditions, varying only the \texttt{dialogue\_prob} and \texttt{battle\_prob} parameters.
All LLM calls are handled by \texttt{PaperFakeLlmClient}, which returns diverse, pre-categorized plan templates to eliminate model variance.
Results are printed as a comparison table and saved to \texttt{backend/data/paper\_experiment\_results.json}.

\section{System Statistics}

As of submission, the Digimon World system includes:
\begin{itemize}[leftmargin=*,nosep]
    \item 860+ passing tests across all subsystems.
    \item 30 species of Digimon agents.
    \item 6 world regions with 14+ sub-regions.
    \item 5-stage evolutionary system with special evolution pathways.
    \item Type-advantage battle system (Vaccine/Data/Virus/Free).
    \item MBTI personality engine with 4-dimensional continuous profiles.
    \item Ebbinghaus-based memory system with importance-weighted retrieval.
    \item 5-tier relational circle model (差序格局).
    \item CPM-grounded continuous mood evolution.
    \item Cooperation threshold system for alliance formation.
    \item Economy, weather, day/night, ecology, disaster, festival, and faction subsystems.
\end{itemize}

\end{document}
